\documentclass[aspectratio=169,11pt]{beamer}

\usepackage[T1]{fontenc}
% `provide=*` permite usar los datos modernos de Babel incluso en una
% instalación compacta que no incluya spanish.ldf.
\usepackage[provide=*,spanish]{babel}

\usetheme{Teaching}

\title[Álgebra lineal]{Álgebra lineal y geometría I}
\subtitle{Una muestra del tema Teaching}
\author[F. Muro]{Fernando Muro}
\institute[US]{Universidad de Sevilla}
\date{Curso 2026--2027}

\begin{document}

\begin{frame}[plain]
  \titlepage
\end{frame}

\begin{frame}{Índice}
  \tableofcontents
\end{frame}

\section{Aplicaciones lineales}

\begin{frame}{Núcleo e imagen}
  \begin{definition}
    Sea $f\colon V\to W$ una aplicación lineal. Definimos
    \[
      \ker f=\{v\in V:f(v)=0\},
      \qquad
      \operatorname{im}f=\{f(v):v\in V\}.
    \]
  \end{definition}

  \begin{proposition}
    Tanto $\ker f$ como $\operatorname{im}f$ son subespacios vectoriales.
  \end{proposition}
\end{frame}

\begin{frame}{Teorema de la dimensión}
  \begin{theorem}
    Si $V$ es de dimensión finita, entonces
    \[
      \dim V=\dim(\ker f)+\dim(\operatorname{im}f).
    \]
  \end{theorem}

  \begin{proof}
    Completamos una base de $\ker f$ hasta obtener una base de $V$ y
    aplicamos $f$ a los vectores añadidos.
  \end{proof}
\end{frame}

\begin{frame}{Consecuencias y ejemplos}
  \begin{corollary}
    Si $\dim V=\dim W<\infty$, entonces $f$ es inyectiva si y solo si es
    suprayectiva.
  \end{corollary}

  \begin{example}
    La aplicación $f\colon\mathbb{R}^3\to\mathbb{R}^2$ dada por
    \[
      f(x,y,z)=(x+y,y+z)
    \]
    tiene núcleo generado por $(1,-1,1)$.
  \end{example}
\end{frame}

\begin{frame}{Una lista progresiva}
  Para estudiar una aplicación lineal:
  \begin{enumerate}
    \item elegimos bases;
    \pause
    \item calculamos su matriz;
    \pause
    \item reducimos mediante eliminación de Gauss.
  \end{enumerate}

  \medskip
  \alert{Idea clave:} la matriz depende de las bases; la aplicación no.
\end{frame}

\end{document}
